\section{Ekspresi Reguler: Definisi dan Operasi}
\label{sec:ch05-regex}

\subsection{Definisi}

\begin{definisi}[Ekspresi reguler]
Diberikan alfabet \(\Sigma\). Ekspresi reguler atas \(\Sigma\)
didefinisikan secara rekursif:
\begin{enumerate}[label=(\roman*),leftmargin=*]
  \item \(\varnothing\), \(\varepsilon\), dan setiap \(a \in \Sigma\)
        adalah ekspresi reguler;
  \item jika \(r\) dan \(s\) ekspresi reguler, maka \(r \mid s\),
        \(rs\), dan \(r^{*}\) juga ekspresi reguler;
  \item tidak ada bentuk lain.
\end{enumerate}
Bahasa yang dinyatakan \(r\), ditulis \(L(r)\), didefinisikan
sebagai himpunan string yang cocok dengan \(r\).
\end{definisi}

Operator utama adalah gabungan (\eng{union}) \(r \mid s\),
konkatenasi \(rs\), dan bintang Kleene \(r^{*}\). Operator
\(r^{+}\) adalah singkatan \(rr^{*}\), dan \(r?\) adalah
\(r \mid \varepsilon\).

\subsection{Prioritas Operator}

Prioritas operator dari tertinggi ke terendah adalah bintang Kleene,
konkatenasi, lalu gabungan. Dengan demikian
\(ab^{*} \mid c\) dibaca \((a(b^{*})) \mid c\), bukan
\(a(b^{*} \mid c)\) atau \((ab)^{*} \mid c\).

\begin{figure}[H]
\centering
\begin{tikzpicture}[node distance=5mm and 7mm, scale=0.9, transform shape]
  \node[phase, text width=1.6cm] (a) {\(a\)};
  \node[phase, text width=1.6cm, right=of a] (b) {\(b\)};
  \node[phase, text width=2.0cm, right=of b] (u) {\(a \mid b\)};
  \node[phase, text width=2.4cm, right=of u] (s) {\((a \mid b)^{*}\)};
  \node[phase, text width=3.0cm, right=of s] (c) {\((a \mid b)^{*}abb\)};
  \draw[arrow] (a) -- (u);
  \draw[arrow] (b) -- (u);
  \draw[arrow] (u) -- (s);
  \draw[arrow] (s) -- (c);
\end{tikzpicture}
\caption{Komposisi ekspresi reguler \((a \mid b)^{*}abb\)}
\label{fig:ch05-komposisi}
\end{figure}

\subsection{Identitas Aljabar}

Ekspresi reguler memenuhi sejumlah identitas yang berguna untuk
penyederhanaan:
\[
\begin{aligned}
r \mid s &= s \mid r, &
(rs)t &= r(st),\\
r \mid r &= r, &
r\varepsilon &= \varepsilon r = r,\\
\varnothing \mid r &= r, &
\varnothing r &= r\varnothing = \varnothing,\\
(r \mid s)t &= rt \mid st, &
r(s \mid t) &= rs \mid rt,\\
(r^{*})^{*} &= r^{*}, &
\varepsilon^{*} &= \varepsilon,\\
\varnothing^{*} &= \varepsilon, &
(r \mid \varepsilon)^{*} &= r^{*}.
\end{aligned}
\]

Identitas ini dipakai oleh alat seperti Flex untuk menyederhanakan
otomata dan mempercepat pencocokan. Namun, identitas tidak selalu
mempermudah pembacaan manusia. Untuk spesifikasi token, kejelasan
lebih penting daripada bentuk paling ringkas.

\subsection{Ekspresi Reguler untuk Token}

Berikut pola umum token:
\begin{itemize}[leftmargin=*]
  \item identifier: \verb![A-Za-z_][A-Za-z0-9_]*!
  \item integer: \verb![0-9]+!
  \item float: \verb![0-9]+\.[0-9]+!
  \item string: \verb!"([^"\\]|\\.)*"!
  \item komentar baris: \verb!//[^\n]*!
  \item komentar blok: \verb!/\*([^*]|\*[^/])*\*/!
\end{itemize}

Notasi di atas memakai kelas karakter dan escape. Dalam implementasi,
detail seperti escape sequence dan karakter Unicode perlu
diperhatikan. Untuk Subset-C, kita membatasi diri pada ASCII agar
spesifikasi dan pengujian lebih sederhana.

\begin{warningbox}
Ekspresi reguler untuk komentar blok di atas adalah penyederhanaan.
Implementasi harus menangani kasus seperti \code{/**/} dan
komentar yang tidak ditutup. Uji selalu dengan kasus batas.
\end{warningbox}
